\section{方法}

\subsection{方法总览}

本文方法记为 EA-RG-MAPPO-S，即 Edge-Aware Role Graph MAPPO with Staged random-radius fine-tuning。方法框图如图~\ref{fig:method_overview} 所示。EA-RG-MAPPO-S 由局部观测编码器、边特征增强角色图编码器、集中式 critic 和分阶段随机通信半径微调组成。策略输入由局部观测表示和图表示拼接得到，critic 使用集中式全局状态。

\begin{figure*}[t]
    \centering
    \includegraphics[width=0.92\textwidth]{method_overview_ea_rg_mappo_s.png}
    \caption{EA-RG-MAPPO-S 方法框图。}
    \label{fig:method_overview}
\end{figure*}

\subsection{角色图构建}

每个时间步构建动态图 $G^t=(V^t,E^t)$。节点集合包括追击无人机节点和目标节点。每个节点包含归一化位置、航向、速度等运动状态，同时使用角色标记区分 UAV 和 target。节点特征与 role embedding 拼接后输入图编码器：
\begin{equation}
h_i^0 = f_{\mathrm{in}}\left([x_i,\mathrm{Emb}(\mathrm{role}_i)]\right).
\end{equation}

角色图使模型能够区分队友节点和目标节点，避免将所有节点视为同质实体。

\subsection{相对边特征增强注意力}

对节点 $i$ 到节点 $j$ 的边，构造相对边特征 $e_{ij}$，包含相对位置、归一化距离、相对通信半径距离、方位角、相对速度、通信可达标记和目标节点标记：
\begin{equation}
e_{ij}=[\Delta x_{ij},\Delta y_{ij},d_{ij}/d_{\max},d_{ij}/R_c,
\cos\beta_{ij},\sin\beta_{ij},\Delta v^x_{ij},\Delta v^y_{ij},
\mathbb{I}_{ij}^{\mathrm{comm}},\mathbb{I}_{j}^{\mathrm{target}}].
\end{equation}
其中 $d_{ij}$ 为节点间距离，$d_{\max}$ 为环境尺度归一化常数，$\beta_{ij}$ 为相对方位角，$\mathbb{I}_{ij}^{\mathrm{comm}}$ 表示通信半径内可达关系。边特征增强图注意力得分为：
\begin{equation}
\mathrm{score}_{ij}=\mathrm{LeakyReLU}\left(a^\top[Wh_i,Wh_j]\right)+g_e(e_{ij}),
\end{equation}
其中 $g_e$ 为边特征打分网络。随后对可通信邻居进行 masked softmax：
\begin{equation}
\alpha_{ij}=\frac{\exp(\mathrm{score}_{ij})}{\sum_{k\in\mathcal{N}_i}\exp(\mathrm{score}_{ik})}, \quad j\in\mathcal{N}_i.
\end{equation}
节点表示更新为：
\begin{equation}
h_i'=\tanh\left(\sum_{j\in\mathcal{N}_i}\alpha_{ij}Wh_j\right).
\end{equation}

这种设计使注意力机制同时考虑节点语义和边的物理关系。与仅由节点特征计算注意力的普通 GAT 不同，EA-RG-MAPPO-S 可以直接利用距离、方位、速度差和通信可达性调整邻居权重。注意力热力图显示，通信半径较小时，部分队友边不可达，UAV 节点主要关注自身和目标；通信半径较大时，注意力在队友和目标之间分布更均匀。

\subsection{MAPPO 优化}

策略网络输出每个智能体的动作分布 $\pi_\theta(a_i|o_i,G)$。采用 MAPPO clipped objective。概率比定义为：
\begin{equation}
\rho_i^t(\theta)=\frac{\pi_\theta(a_i^t|o_i^t,G^t)}{\pi_{\theta_{\mathrm{old}}}(a_i^t|o_i^t,G^t)}.
\end{equation}

策略损失为：
\begin{equation}
L_{\mathrm{policy}}=-\mathbb{E}\left[\min\left(\rho_i^t A_i^t,\mathrm{clip}(\rho_i^t,1-\epsilon,1+\epsilon)A_i^t\right)\right].
\end{equation}

总损失为：
\begin{equation}
L=L_{\mathrm{policy}}+c_vL_{\mathrm{value}}-c_HL_{\mathrm{entropy}}+c_{\mathrm{aux}}L_{\mathrm{aux}}.
\end{equation}
其中 $L_{\mathrm{aux}}$ 是可选辅助损失。由于当前目标意图分支的 balanced accuracy 不足，本文主结论不依赖该辅助分支。

\subsection{分阶段随机通信半径微调}

固定通信半径训练容易使策略适应单一通信拓扑。本文采用两阶段训练。第一阶段在固定通信半径 $R_c=8$ 下训练 edge-aware role graph policy，使模型学到基础协同能力。第二阶段从第一阶段 checkpoint 出发，以较小学习率进行短程微调。每个 episode 从区间中随机采样通信半径：
\begin{equation}
R_c\sim U(4,10).
\end{equation}

该阶段的目标是增强策略对不同通信拓扑的适应性，而不是简单增加训练时间。
